(1) 外场为纯偶极角依赖，球内在原点有限的解只能写 \(\Phi_{\rm in}=-A E_0r\cos\theta\)，外部满足无穷远均匀场的解写成 \(\Phi_{\rm out}=-E_0r\cos\theta+C E_0a^3\cos\theta/r^2\)。无净自由电荷故无需外部 \(1/r\) 单极项。
(2) 在 \(r=a\)，势连续给 \(A=1-C\)；法向位移连续 \(\varepsilon_1E_{r,\rm in}=\varepsilon_2E_{r,\rm out}\)，即 \(\varepsilon_1 A=\varepsilon_2(1+2C)\)。联立得 \(A=3\varepsilon_2/(\varepsilon_1+2\varepsilon_2)\)、\(C=(\varepsilon_1-\varepsilon_2)/(\varepsilon_1+2\varepsilon_2)\)；内部均匀电场 \(\mathbf E_{\rm in}=AE_0\hat z\)。
(3) 外部偶极势为 \(p\cos\theta/(4\pi\varepsilon_2r^2)\)，因此 \(p=4\pi\varepsilon_2a^3 C E_0\)，即 \(\alpha_{\rm pol}=p/E_0=4\pi\varepsilon_2a^3(\varepsilon_1-\varepsilon_2)/(\varepsilon_1+2\varepsilon_2)\)。当球介电常数较大时 \(C>0\)，诱导偶极与外场同向。
(4) \(\varepsilon_1=\varepsilon_2\) 时 \(C=0,A=1\)，球不改变场；\(\varepsilon_1\to\infty\) 时 \(A\to0,C\to1\)，恢复中性理想导体球的内零场与偶极矩 \(4\pi\varepsilon_2a^3E_0\)。有限正介电常数时 \(A\) 一般非零，不能未经计算就套用导体边界条件。
